% ECONOMICS_PROBLEM: {"id":"C73-11","title":"Characterizing universal convergence of fictitious play","jel_code":"C73","source_id":"OP-0136"}
% FIRST_POSTED: 2026-10-10
% ATTEMPT_STATUS: solved
% ENTRY_KIND: research_attempt
% PROOF_SCOPE: Complete literal payoff-array equivalence; explicitly infinitary.
% !TeX program = pdflatex
\documentclass[11pt,a4paper]{article}
\usepackage[T1]{fontenc}
\usepackage{lmodern}
\usepackage[left=20mm,right=20mm,top=18mm,bottom=18mm]{geometry}
\usepackage{amsmath,amssymb,amsthm,mathtools}
\usepackage{microtype,enumitem,xurl,hyperref}
\hypersetup{hypertexnames=false,colorlinks=true,linkcolor=blue,citecolor=blue,urlcolor=blue,
  pdftitle={C73-11: An infinitary payoff-array characterization},
  pdfauthor={Ji Ho Bae}}
\setlength{\parindent}{0pt}
\setlength{\parskip}{3pt}
\setlist[enumerate]{itemsep=2pt,topsep=4pt}
\newtheorem{theorem}{Theorem}
\newtheorem{lemma}[theorem]{Lemma}
\newtheorem{proposition}[theorem]{Proposition}
\theoremstyle{definition}
\newtheorem{definition}[theorem]{Definition}
\newtheorem{example}[theorem]{Example}
\newcommand{\R}{\mathbb R}
\newcommand{\N}{\mathbb N}
\newcommand{\CFP}{\mathcal C_{\mathrm{FP}}}
\DeclareMathOperator{\NE}{NE}
\DeclareMathOperator{\dist}{dist}
\DeclareMathOperator{\rank}{rank}
\DeclareMathOperator*{\argmax}{arg\,max}
\title{C73-11: An infinitary payoff-array characterization\newline
  of universal convergence of fictitious play}
\author{Ji Ho Bae}
\date{10 October 2026}

\begin{document}
\maketitle
\section*{Model, process, and claim status}
\noindent\textbf{Model:} GPT Codex. \textbf{Contributor:} Ji Ho Bae.
The exact serving revision was not exposed.
The complete mathematical argument below was generated by GPT Codex,
using parallel derivation and adversarial proof checking in response to a
request for a complete hand proof of the exact catalogue statement.
No human-written mathematical edits are included. The original theorem,
definitions, proofs, verification, citations, and their order are
preserved in full.
This is an LLM claim of a complete payoff-array characterization of
universal convergence of fictitious play; the criterion is explicitly
infinitary, and community expert review is pending.
No proof-assistant verification is claimed.
The complete original generated output is preserved at the immutable
\href{https://github.com/jbaelaw/universal-fictitious-play-characterization/blob/09c4bb2c4b9b7a64c2a9fb87516da03a1b559abc/paper/Proof.tex}{manuscript TeX}
and \href{https://github.com/jbaelaw/universal-fictitious-play-characterization/blob/09c4bb2c4b9b7a64c2a9fb87516da03a1b559abc/paper/Proof.pdf}{manuscript PDF}.

\begin{abstract}
For every finite normal-form game we construct a countable tree from
finite polynomial-sign tests on its payoff entries. Simultaneous fictitious
play approaches the Nash set from every initial profile and under every
legal tie choice if and only if this tree admits a strictly decreasing
countable-ordinal labeling. The proof retains a common initial profile
through nested compact feasible sets. The criterion is explicitly
infinitary. It also yields a coanalytic upper bound for the convergence
class and an elementary example whose tree has root rank $\omega$.
\end{abstract}

\section{The exact problem and the claim}

Let $G=(N,(S_i),(u_i))$ be a finite normal-form game, where
$N=\{1,\ldots,n\}$, $n\ge1$, and every $S_i$ is nonempty. Put
$A=\prod_iS_i$ and $X=\prod_i\Delta(S_i)$. Payoff entries are arbitrary
real numbers; payoffs on $X$ are their multilinear extensions. Starting
from any $x^1\in X$, the specified simultaneous fictitious-play rule is
\begin{equation}\label{eq:fp}
 \begin{aligned}
 a_i^{t+1}&\in\argmax_{b\in S_i}u_i(b,x_{-i}^t),\\
 x_i^{t+1}&=\frac{t x_i^t+e_{a_i^{t+1}}}{t+1}\quad(t\ge1).
 \end{aligned}
\end{equation}
Every legal, time-dependent selection among ties is permitted. The
catalogue problem C73-11 asks for necessary and sufficient conditions on
the finite payoff arrays for membership in
\begin{equation}\label{eq:class}
 \CFP=\left\{G:
 \begin{array}{l}
 % Browser-equivalent formatting only; mathematical content is identical.
% The immutable original standalone manuscript and PDF remain preserved.
\text{for every }x^1\in X\text{ and every legal sequence in (1),}\\[-2pt]
 \displaystyle\dist(x^t,\NE(G))\longrightarrow0
 \end{array}\right\},
 \qquad
 \dist(x,K)=\inf_{z\in K}\|x-z\|_2.
\end{equation}

The condition below is a necessary and sufficient condition for this
entire class. It consists of a countable tree determined by finite
quantifier-free payoff tests, together with an ordinal-ranking condition.
It is not a finite membership algorithm or a classification by familiar
game classes. The learning-dynamics discussion and theoretical open item~5
of Li et al.~\cite{li} motivate the source agenda; the precise rule and
quantifiers in \eqref{eq:class} are the editorial C73-11 specialization.
No broader classification agenda is claimed resolved.

\section{Explicit finite polynomial tests}

Introduce the finite real vector $\zeta=(\zeta_{i,b})$ and the simplex
constraints
\begin{equation}\tag{S}\label{eq:S}
 \zeta_{i,b}\ge0,\qquad \sum_{b\in S_i}\zeta_{i,b}=1
 \quad(i\in N,\ b\in S_i).
\end{equation}
Fix an integer $T\ge1$ and a finite joint-action word
$v=(a^2,\ldots,a^T)$, with $a^q\in A$. The word is empty when $T=1$.
For $1\le t\le T$, define
\begin{align}
 c_{i,b}^t(v)&=\#\{q\in\{2,\ldots,t\}:a_i^q=b\},
 &Z_{i,b}^t(\zeta)&=\zeta_{i,b}+c_{i,b}^t(v),\label{eq:counts}\\
 P_{i,b}^t(\zeta;u)&=
 \sum_{s_{-i}\in\prod_{j\ne i}S_j}
 u_i(b,s_{-i})\prod_{j\ne i}Z_{j,s_j}^t(\zeta),\label{eq:P}\\
 H_{i,b}^t(\zeta;u)&=
 tP_{i,b}^t(\zeta;u)
 -\sum_{d\in S_i}Z_{i,d}^t(\zeta)P_{i,d}^t(\zeta;u).
 \label{eq:H}
\end{align}
For $n=1$, the empty product in \eqref{eq:P} is one and the sum has one
term. These are explicit polynomials in $\zeta$ and the payoff entries,
with rational coefficients determined by the finite integer/action data.
Impose the legality inequalities
\begin{equation}\tag{L}\label{eq:L}
 P_{i,a_i^{t+1}}^t-P_{i,b}^t\ge0
 \quad(1\le t<T,\ i\in N,\ b\in S_i).
\end{equation}

A \emph{record} consists of this word, a positive integer $m$, marked
times $1\le t_1<\cdots<t_r=T$, $r\ge1$, and witnesses
$h_\ell=(i_\ell,b_\ell)$ with $b_\ell\in S_{i_\ell}$. Add
\begin{equation}\tag{B}\label{eq:B}
 mH_{i_\ell,b_\ell}^{t_\ell}-t_\ell^n\ge0
 \quad(1\le\ell\le r).
\end{equation}
Write $C_q(u)\subseteq X$ for the $\zeta$ satisfying
\eqref{eq:S}, \eqref{eq:L}, and \eqref{eq:B} for record $q$.
It is closed and compact, and it is semialgebraic. A record is
\emph{admitted} precisely when $C_q(u)\ne\varnothing$.

For a definition entirely in the payoff entries, apply real quantifier
elimination~\cite{tarski} to the finite formula
\begin{equation}\label{eq:QE}
 \exists\zeta\;\bigl[% Browser-equivalent formatting only; mathematical content is identical.
% The immutable original standalone manuscript and PDF remain preserved.
\text{(S), (L), and (B)}
 \bigr],
\end{equation}
treating every $u_i(s)$ as a free variable. Denote a resulting
quantifier-free finite Boolean combination of polynomial signs by
$\Theta_q(u)$. Thus
\begin{equation}\label{eq:admission}
 C_q(u)\ne\varnothing\quad\Longleftrightarrow\quad\Theta_q(u).
\end{equation}
The compilation depends only on $q$ and eliminates $\zeta$ entirely.
It works for arbitrary real arrays. For rational arrays, each individual
finite test is decidable. This does not decide the later well-foundedness
condition on the countable tree.

\section{The payoff-generated tree and the main theorem}

Define $\mathcal T(u)$ to have a root $e$ and one nonroot vertex for each
record satisfying $\Theta_q(u)$. A record with $r=1$ marks has parent $e$.
For $r\ge2$, its parent deletes the last mark and witness and truncates
the word after $t_{r-1}$. Extension always preserves $m$ and every earlier
action, mark, and witness. Equivalently, an edge appends a nonempty finite
action block ending at a new marked time; only the first block can be empty,
when the first mark is at time one.

Each admitted record has an admitted parent, since its feasible set is
contained in the parent's feasible set. The root is always present.
There are countably many possible records. The tree may have infinitely
many children at a vertex because the next marked time is unbounded.
Its definition by \eqref{eq:admission} uses only payoff entries and finite
integer/action labels.

\begin{definition}[Payoff-array condition (O)]\label{def:O}
There exist a countable ordinal $\alpha$ and a labeling
$\rho:\mathcal T(u)\to\alpha$ such that
\begin{equation}\tag{O}\label{eq:O}
 \rho(w)<\rho(q)\qquad\text{whenever $w$ is a child of $q$}.
\end{equation}
The labeling includes the root. Neither a finite encoding of $\alpha$
nor computability of $\rho$ is required.
\end{definition}

Condition (O) contains no trajectory, initial-profile variable, Nash-set
variable, limit, or original convergence assertion: its admitted vertices
are selected by the compiled payoff formulas $\Theta_q$.

\begin{theorem}\label{thm:main}
For every finite normal-form game as above, the following are equivalent:
\begin{enumerate}[label=\textup{(\roman*)}]
\item $G\in\CFP$, with all the quantifiers and the simultaneous rule in
\eqref{eq:class}.
\item Its payoff array satisfies condition \textup{(O)}.
\item $\mathcal T(u)$ is well founded, meaning that it has no infinite branch.
\end{enumerate}
\end{theorem}

\section{Exact interpretation and the Nash residual}

\begin{lemma}\label{lem:finite}
For an initial profile $\zeta\in X$ and a finite word $v$, the update
\eqref{eq:fp} has the identities
\begin{equation}\label{eq:identity}
 \begin{aligned}
 x_{i,b}^t&=\frac{Z_{i,b}^t}{t},\qquad
 u_i(b,x_{-i}^t)=\frac{P_{i,b}^t}{t^{n-1}},\\
 u_i(b,x_{-i}^t)-u_i(x^t)&=\frac{H_{i,b}^t}{t^n}.
 \end{aligned}
\end{equation}
Consequently \textup{(L)} is exactly legality of the finite word, and
\textup{(B)} is exactly the specified gain of at least $1/m$ at each mark.
\end{lemma}

\begin{proof}
The first identity holds at $t=1$. Multiplying the update by $t+1$ adds
one to exactly the next chosen action's count, proving it by induction.
Substituting it in the multilinear payoff gives the second identity.
Since
$u_i(x^t)=\sum_d(Z_{i,d}^t/t)(P_{i,d}^t/t^{n-1})$,
the third follows from \eqref{eq:H}. All denominators are positive, so
clearing them gives precisely \eqref{eq:L} and \eqref{eq:B}, including
weak ties.
\end{proof}

Define
\begin{equation}\label{eq:residual}
 R(x)=\max_{i\in N,\,b\in S_i}
 \bigl[u_i(b,x_{-i})-u_i(x)\bigr].
\end{equation}
For each player the $x_i$-weighted mean of the bracketed gains is zero,
so $R\ge0$. It is continuous, and $R(x)=0$ exactly when no pure deviation
improves payoff. A mixed deviation is a convex combination of pure ones;
thus the zero set is exactly $\NE(G)$.

The zero set is nonempty by Nash's finite-game existence theorem~\cite{nash}.
For completeness, the following continuous-map argument proves existence
using Brouwer's fixed-point theorem.
Set $g_{i,b}(x)=[u_i(b,x_{-i})-u_i(x)]_+$,
$G_i(x)=\sum_b g_{i,b}(x)$, and consider the continuous map
\[
 F_{i,b}(x)=\frac{x_{i,b}+g_{i,b}(x)}{1+G_i(x)}
 \quad\text{from }X\text{ to }X.
\]
Brouwer's fixed-point theorem supplies a fixed point. There,
$g_{i,b}=x_{i,b}G_i$. If $G_i>0$, every action with positive $x_i$-mass
has strictly positive gain, contradicting the zero weighted mean of gains.
Hence all $G_i=0$, giving a Nash equilibrium. Continuity also shows that
the Nash set is compact.

\begin{lemma}\label{lem:distance}
For every sequence $x^t\in X$,
\begin{equation}\label{eq:distance}
 R(x^t)\longrightarrow0
 \quad\Longleftrightarrow\quad
 \dist(x^t,\NE(G))\longrightarrow0.
\end{equation}
Failure is equivalent to the existence of one integer $m\ge1$ and
infinitely many times at which a witnessed pure-deviation gain is at least
$1/m$.
\end{lemma}

\begin{proof}
Uniform continuity of $R$ on compact $X$, its vanishing on the Nash set,
and attainment of distance to that compact set prove the implication from
distance convergence. Conversely, suppose $R(x^t)\to0$ but distance
stays at least some $\varepsilon>0$ along a subsequence. A further
subsequence converges to $x^*\in X$. Then $R(x^*)=0$, so $x^*$ is Nash,
contradicting the positive distance bound.

Since $R\ge0$, failure of its convergence to zero means that
$R(x^t)\ge\varepsilon$ infinitely often for some $\varepsilon>0$.
Choose $m$ with $1/m\le\varepsilon$. The finite maximum in
\eqref{eq:residual} has a witness at every such time. Conversely, infinitely
many gains at least one fixed $1/m$ prevent residual convergence.
\end{proof}

\section{Branches recover one common initial profile}

\begin{proof}[Proof of Theorem~\ref{thm:main}, \textup{(i)}$\Leftrightarrow$\textup{(iii)}]
If a legal fictitious-play orbit fails convergence, Lemma~\ref{lem:distance}
gives one $m$, increasing times $t_1,t_2,\ldots$, and witnesses
$(i_\ell,b_\ell)$. At stage $r$, take the orbit's word through $t_r$
and its first $r$ marks. Its actual initial profile satisfies all the finite
constraints by Lemma~\ref{lem:finite}. These admitted records are successive
vertices of an infinite branch.

Conversely, let $q_1,q_2,\ldots$ be an infinite branch after the root.
The first vertex fixes $m$ forever. Successive marked times strictly
increase and hence tend to infinity. Consistent word extensions specify
an action at every update date. The sets $C_{q_r}(u)$ are nonempty closed
subsets of the same compact $X$, and
\begin{equation}\label{eq:intersection}
 C_{q_{r+1}}(u)\subseteq C_{q_r}(u),\qquad
 \bigcap_{r\ge1}C_{q_r}(u)\ne\varnothing.
\end{equation}
The last assertion is the nested compact-set property. Choose one
$\zeta$ in this intersection. Every legality inequality at any finite
update date occurs in all sufficiently long records, so the union word
is a legal simultaneous orbit from this same $\zeta$. Every marked gain
is at least $1/m$. Its residual therefore fails to converge to zero;
Lemma~\ref{lem:distance} gives failure of Nash-set convergence.
\end{proof}

Retaining all earlier constraints is crucial: separate feasibility of
unrelated records would not recover a common initial profile. Weak
inequalities retain every tie and make \eqref{eq:intersection} applicable.
The fixed threshold is also crucial; a branch cannot increase $m$ to
turn shrinking residuals into an obstruction.

\section{Ordinal rankings and a root rank of \texorpdfstring{$\omega$}{omega}}

\begin{proof}[Proof of Theorem~\ref{thm:main}, \textup{(ii)}$\Leftrightarrow$\textup{(iii)}]
Condition (O) forbids an infinite branch, since ordinals have no infinite
strictly descending sequence. For the converse, let the countable tree be
well founded. Its child relation supports well-founded recursion:
\begin{equation}\label{eq:rank}
 \rank(q)=\sup\{\rank(w)+1:w\text{ is a child of }q\},
 \qquad\sup\varnothing=0.
\end{equation}
By well-founded induction, every rank is countable: assuming this for all
children of a vertex, their ranks have countable supremum because there
are only countably many children. Since there are countably many vertices,
\[
 \alpha=\sup_{q\in\mathcal T(u)}(\rank(q)+1)
\]
is countable. Every rank is less than $\alpha$, and the ranks strictly
decrease along edges, proving (O).

\end{proof}

\begin{example}[An actual fictitious-play tree of root rank $\omega$]
One player has actions $0,1$ with payoffs $0,1$. Action $1$ is the unique
legal update. If the initial mass on action $0$ is $c\in[0,1]$, its mass
at time $t$ is $c/t$, and the Nash residual is $c/t$. Only deviation to
action $1$ can witness a positive gain.

For a record at threshold $1/m$ ending at time $T$, admission requires
$c/T\ge1/m$, hence $T\le m$. Conversely, every increasing list of marked
times ending at $T\le m$ is admitted by $c=1$ and the unique legal word.
An admitted vertex with last time $T$ has at most $m-T$ further marks,
and can realize all times $T+1,\ldots,m$. Its rank is therefore exactly
$m-T$. In particular, the vertex first marked at time one has rank $m-1$.
The combined tree's root has
\[
 \rank(e)=\sup_{m\ge1}\ \sup_{1\le T\le m}(m-T+1)
 =\sup_{m\ge1}m=\omega.
\]
The game converges from every initialization, but its root has no finite
ordinal rank. This makes no claim that higher ranks are realized by this
particular family of dynamics.
\end{example}

\section{A coanalytic payoff-space consequence}

Fix the action sets and identify payoff arrays with $U=\R^D$, where
$D=n|A|$. For each possible record $q$, let
$Q_q=\{u:C_q(u)\ne\varnothing\}$.

\begin{proposition}\label{prop:coanalytic}
Every $Q_q$ is closed and semialgebraic in $U$. The subset $\CFP\subseteq U$
is coanalytic, equivalently a $\Pi^1_1$ set.
\end{proposition}

\begin{proof}
Semialgebraicity follows from \eqref{eq:QE}. To prove closedness, let
$u_k\in Q_q$ converge to $u$ and choose $\zeta_k\in C_q(u_k)$.
Compactness of $X$ supplies a convergent subsequence. All finitely many
polynomial constraints pass to its limit, proving $u\in Q_q$.

Enumerate all syntactically valid complete records by $\N$.
An element $p\in\N^{\N}$ names successive potential vertices
$q(p_0),q(p_1),\ldots$. Define $B\subseteq U\times\N^{\N}$ by requiring:
\begin{enumerate}[label=\textup{(\alph*)}]
\item $q(p_0)$ has one mark;
\item $q(p_{k+1})$ is a child of $q(p_k)$ for every $k$;
\item $u\in Q_{q(p_k)}$ for every $k$.
\end{enumerate}
The first record includes $m$, and the child relation preserves it.
These requirements define a closed set. Explicitly, if
$(u_j,p_j)\to(u,p)$, convergence in Baire space, with its discrete
coordinate topology, makes each fixed finite list of coordinates
eventually constant and equal to the corresponding list of $p$.
Requirements (a) and (b) pass to the limit. For each fixed $k$, the
record named by $(p_j)_k$ is eventually exactly $q(p_k)$, so closedness
of $Q_{q(p_k)}$ passes (c) to the limit as well.

Theorem~\ref{thm:main} says that the projection of $B$ onto $U$ is exactly
$U\setminus\CFP$. A projection of a closed subset of
$U\times\N^{\N}$ is analytic; its complement is coanalytic. This proves
the asserted upper bound.
\end{proof}

This argument establishes no $\Pi^1_1$ hardness or completeness,
non-Borelness, undecidability, or necessity of high ordinal ranks for
fictitious-play games.

\section{Scope and boundary cases}

The initial domain is the entire product of closed simplexes, including
boundary faces and irrational coordinates. Weak best-response inequalities
include all payoff ties. Every player updates simultaneously, and the
target remains distance to the exact Nash set. The proof does not replace
that target by point convergence, action convergence, or convergence of
joint empirical play to a correlated-equilibrium concept.

Singleton action sets cause no change in the formulas. The empty products
for a one-player game have already been specified. If all payoffs are
constant, then $H=0$, no marked inequality is feasible, and the tree has
only its root, of rank zero; every profile is Nash. If a convention admits
an empty-player game, that trivial case may separately be assigned the
root-only tree and is convergent. Positive playerwise payoff rescaling and
adding a function of opponents' actions preserve legal choices and the
Nash set. They may alter a numerical threshold, but not the existence of
some fixed integer obstruction threshold.

\begin{thebibliography}{9}
\bibitem{li}
Hanyu Li, Wenhan Huang, Zhijian Duan, David Henry Mguni, Kun Shao,
Jun Wang, and Xiaotie Deng.
\emph{A survey on algorithms for Nash equilibria in finite normal-form
games}. arXiv:2312.11063v1, 18 December 2023.
Learning-dynamics discussion and concluding ``Open problems and future
directions,'' theoretical item~5, printed page 30.
\href{https://arxiv.org/pdf/2312.11063v1}{Primary v1 PDF}.

\bibitem{nash}
John F. Nash, Jr.
\emph{Equilibrium points in $n$-person games}.
Proceedings of the National Academy of Sciences 36(1), 48--49 (1950).
DOI: 10.1073/pnas.36.1.48. Finite-game existence theorem, proved in the
note by Kakutani. Our displayed Brouwer proof is supplied separately.
\href{https://pmc.ncbi.nlm.nih.gov/articles/PMC1063129/}{Primary paper scan}.

\bibitem{tarski}
Alfred Tarski.
\emph{A Decision Method for Elementary Algebra and Geometry:
Prepared for Publication with the Assistance of J.~C.~C.~McKinsey}.
RAND Corporation, Report R-109, 1951 edition.
Real quantifier elimination is used here only to compile individual finite
formulas, not to decide well-foundedness of a countable tree.
\href{https://www.rand.org/pubs/reports/R109.html}{Edition record}; \href{https://www.rand.org/content/dam/rand/pubs/reports/2008/R109.pdf}{primary scan}.
\end{thebibliography}
\end{document}
